\documentclass[aspectratio=169,11pt]{beamer}
\usepackage[T1]{fontenc}
\usepackage[utf8]{inputenc}
\usepackage[brazilian]{babel}
\usepackage{lmodern,amsmath,amssymb,booktabs,array,tikz}
\usetikzlibrary{arrows.meta,positioning,calc}
\definecolor{Navy}{HTML}{132D46}
\definecolor{Teal}{HTML}{008B8B}
\definecolor{Orange}{HTML}{DA6B34}
\definecolor{Mist}{HTML}{EDF4F7}
\definecolor{Gray}{HTML}{64748B}
\setbeamercolor{normal text}{fg=Navy,bg=white}
\setbeamercolor{structure}{fg=Teal}
\setbeamercolor{frametitle}{fg=Navy}
\setbeamercolor{block title}{fg=white,bg=Navy}
\setbeamercolor{block body}{fg=Navy,bg=Mist}
\setbeamercolor{itemize item}{fg=Teal}
\setbeamertemplate{navigation symbols}{}
\setbeamertemplate{itemize item}{\small\raise1pt\hbox{$\bullet$}}
\setbeamertemplate{itemize subitem}{\tiny$\bullet$}
\setbeamerfont{frametitle}{size=\Large,series=\bfseries}
\setbeamerfont{block title}{size=\normalsize,series=\bfseries}
\setbeamersize{text margin left=8mm,text margin right=8mm}
\setbeamertemplate{frametitle}{\vspace{2mm}{\scriptsize\color{Teal}\insertsectionhead\par}\vspace{1mm}\insertframetitle\par\vspace{2mm}}
\newcommand{\MainTotal}{32}
\newif\ifbackup
\setbeamertemplate{footline}{\hspace{8mm}\begin{minipage}{.80\paperwidth}\color{Gray}\tiny Gráviton de Souza\enspace /\enspace Defesa de mestrado\enspace /\enspace ITA\end{minipage}\hfill{\color{Teal}\scriptsize\ifbackup Apoio\else\insertframenumber/\MainTotal\fi}\hspace{7mm}\vspace{3mm}}
\newcommand{\take}[1]{\par\vspace{2mm}\begin{beamercolorbox}[wd=\linewidth,sep=2.5mm]{block body}\small\textbf{#1}\end{beamercolorbox}}
\newcommand{\source}[1]{\par\vspace{1mm}{\color{Gray}\fontsize{7.5}{9}\selectfont #1\par}}
\newcommand{\fig}[2][.95]{\centering\includegraphics[width=#1\linewidth,height=.49\textheight,keepaspectratio]{#2}\par}
\newcommand{\bigword}[2]{\textcolor{Teal}{\fontsize{30}{34}\selectfont\bfseries #1}\par\vspace{1mm}{\small #2}}
\newcommand{\CN}{\mathcal{CN}}
\newcommand{\KL}{D_{\mathrm{KL}}}
\title{Massa do gráviton e polarizações de ondas gravitacionais em PTAs}
\subtitle{Compressão em frequência e validação estatística}
\author{Gráviton de Souza}
\institute{Instituto Tecnológico de Aeronáutica\\Programa de Pós-Graduação em Física · Física Nuclear}
\date{}
\begin{document}
\setlength{\abovedisplayskip}{6pt}
\setlength{\belowdisplayskip}{6pt}
\section{A pergunta física}
% 01
\begin{frame}[plain]
\begin{tikzpicture}[remember picture,overlay]
\fill[Navy] (current page.south west) rectangle (current page.north east);
\fill[Teal] ([xshift=-38mm]current page.south east) rectangle (current page.north east);
\foreach \r in {11,19,27,35,43}{\draw[white,opacity=.18,line width=1pt] ([xshift=-19mm,yshift=42mm]current page.south east) circle (\r mm);}
\fill[Orange] ([xshift=-19mm,yshift=42mm]current page.south east) circle (2mm);
\end{tikzpicture}
\begin{minipage}{.75\textwidth}
\color{white}\small ITA / FÍSICA / DEFESA DE MESTRADO\par\vspace{6mm}
{\raggedright\fontsize{23}{27}\selectfont\bfseries Massa do gráviton\\e polarizações de ondas\\gravitacionais em PTAs\par}
\vspace{4mm}{\large Compressão em frequência e validação estatística\par}
\vspace{7mm}\textbf{Gráviton de Souza}\par\vspace{1mm}{\small Orientador: Prof. Dr. Schrödinger GPT de Gotham}
\end{minipage}
\end{frame}
% 02
\begin{frame}{O que um limite de massa realmente mede?}
\begin{columns}[T]
\column{.57\textwidth}\vspace{3mm}{\Large Uma posterior estreita pode refletir\par\vspace{3mm}\textcolor{Teal}{informação dos dados},\par\vspace{2mm}a \textcolor{Orange}{aproximação adotada}\par\vspace{2mm}ou o \textbf{suporte da priori}.}
\column{.38\textwidth}\begin{block}{Pergunta da pesquisa}Como separar perda de informação por compressão, erro de verossimilhança e erro de resposta dispersiva?\end{block}
\end{columns}
\end{frame}
% 03
\begin{frame}{Três passos para responder à pergunta}
\vspace{4mm}\begin{columns}[T]
\column{.30\textwidth}\bigword{01}{\textbf{Construir a resposta}}\vspace{3mm}\par Propagação massiva, temporização e polarizações.
\column{.30\textwidth}\bigword{02}{\textbf{Separar as aproximações}}\vspace{3mm}\par Dados pareados e controle gaussiano.
\column{.30\textwidth}\bigword{03}{\textbf{Medir o que se aprende}}\vspace{3mm}\par Calibração, informação e dependência da priori.
\end{columns}\vspace{7mm}\take{A defesa conecta a física da resposta à interpretação estatística dos resultados.}
\end{frame}
% 04
\begin{frame}{Um PTA compara relógios distribuídos pelo céu}
\begin{columns}
\column{.56\textwidth}
\centering\begin{tikzpicture}[scale=.85,>=Latex]
\fill[Mist] (0,0) circle(2.2);\fill[Teal] (0,0) circle(.19);\node[below=3mm] at (0,0){Terra};
\foreach \x/\y/\lab in {-2.6/1.7/1,2.7/1.5/2,1.8/-1.6/3}{\fill[Navy] (\x,\y) circle(.09);\draw[Orange,line width=1.3pt,->] (\x,\y)--(.15,0);\node[above] at (\x,\y){\small Pulsar \lab};}
\foreach \x in {-2.7,-1.9,-1.1,-.3,.5,1.3,2.1}{\draw[Teal,opacity=.28,line width=1pt] (\x,2.5)--(\x+.7,-2.2);}
\node[Teal] at (0,2.9){Frente de onda gravitacional};
\end{tikzpicture}
\column{.41\textwidth}\begin{itemize}\item Pulsos fornecem tempos de chegada.\item O modelo de timing prevê sua evolução.\item Ondas gravitacionais deixam resíduos correlacionados.\end{itemize}
\take{A assinatura está na correlação entre direções, frequências e fases.}
\end{columns}\source{Ilustração esquemática. Hellings--Downs (1983); Liang--Trodden (2021).}
\end{frame}
% 05
\begin{frame}{A massa muda a velocidade de propagação}
\begin{columns}
\column{.58\textwidth}\fig{figures/dispersao.pdf}
\column{.38\textwidth}\[f_g=\frac{m_gc^2}{h_{\rm P}}\]\[u=f_gT,\quad f_k=\frac{k}{T}\]\[\beta_k=\frac{v_g}{c}=\sqrt{1-(u/k)^2}\]\small O primeiro canal está mais próximo do limiar de propagação.
\end{columns}\take{Uma única frequência de referência pode alterar a resposta de cada canal.}\source{Curvas analíticas, não dados simulados. Frequências em Hz; $u$ e $\beta$ adimensionais.}
\end{frame}
% 06
\begin{frame}{A resposta combina geometria e fase do pulsar}
\[R_a^A=\underbrace{\frac{p_a^ip_a^j e^A_{ij}}{2(1+\beta\widehat\Omega\cdot p_a)}}_{\text{geometria e dispersão}}\;
\underbrace{\left[1-e^{-iy_a(1+\beta\widehat\Omega\cdot p_a)}\right]}_{\text{Terra e pulsar}},\qquad y_a=\frac{2\pi f L_a}{c}.\]
\vspace{3mm}\[\Gamma_{ab}(f)=\frac{3}{8\pi}\int d\Omega\sum_{A=+,\times}R_a^A R_b^{A*}.\]
\take{A ORF, função de redução de sobreposição, transporta essa física para a covariância dos resíduos.}
\source{Convenção tensorial de norma dois. Liang--Trodden (2021, versão corrigida); Cordes et al. (2025).}
\end{frame}
% 07
\begin{frame}{Do TG pioneiro à teoria de campo de Fierz--Pauli}
\begin{columns}[T]
\column{.47\textwidth}\begin{block}{Ponto de partida no ITA (2003)}Resgate do estudo pioneiro de Wayne de Paula.\par\vspace{3mm}Curvatura, tetradas e polarizações lineares sob a formulação histórica de Visser.\par\vspace{3mm}Relevância da faixa de baixas frequências.\end{block}
\column{.47\textwidth}\begin{block}{Atualização: Fierz--Pauli linear}Teoria consistente de campo massivo de spin-2.\par\vspace{2mm}Vínculos covariantes eliminam modos espúrios:\[H=0,\qquad q^\mu H_{\mu\nu}=0.\]Restam os 5 graus de liberdade físicos.\end{block}
\end{columns}\take{A dissertação conecta o cálculo histórico à resposta moderna e consistente de PTAs.}\source{TG de Wayne L. S. de Paula (2003); de Paula, Miranda e Marinho (2004); Liang--Trodden (2021).}
\end{frame}
% 08
\begin{frame}{A helicidade zero tem duas deformações vinculadas}
\begin{columns}
\column{.48\textwidth}\centering\begin{tikzpicture}[scale=.9]
\draw[Gray,dashed] (0,0) circle(1);\draw[Teal,line width=2pt] (0,0) ellipse (1.45 and 1.45);
\foreach \a in {0,45,...,315}{\fill[Teal] (\a:1.45) circle(.055);}\node at (0,-2){Transversal: breathing};
\end{tikzpicture}
\column{.48\textwidth}\[\boxed{\frac{p_l}{p_b}=-\left(\frac{f_g}{f}\right)^2}\]\small Na convenção de projeções de curvatura utilizada.\par\vspace{4mm}Uma única amplitude escalar produz componentes transversal e longitudinal.
\end{columns}\take{A análise respeita o vínculo físico; não ajusta dois modos escalares independentes.}\source{Fierz--Pauli linear; conversão entre bases apresentada no material de apoio.}
\end{frame}
\section{Desenho da comparação}
% 09
\begin{frame}{Três mudanças que precisam ser distinguidas}
\centering\begin{tikzpicture}[>=Latex,node distance=9mm,box/.style={draw=Teal,fill=Mist,rounded corners=2mm,minimum width=29mm,minimum height=13mm,align=center}]
\node[box](a0){$A_0$\\Coeficientes};\node[box,right=of a0](a){$A$\\Estimadores};\node[box,right=of a](b){$B$\\Compressão};\node[box,right=of b](c){$C$\\Resposta fixa};
\draw[->,thick](a0)--(a);\draw[->,thick](a)--(b);\draw[->,thick](b)--(c);
\node[below=6mm of a,align=center,font=\small]{Distribuição\\aproximada};\node[below=6mm of b,align=center,font=\small]{Menos\\coordenadas};\node[below=6mm of c,align=center,font=\small]{Outra resposta\\em frequência};
\end{tikzpicture}\vspace{6mm}\take{A mesma realização é usada nos contrastes: muda-se uma escolha identificável por vez.}
\end{frame}
% 10
\begin{frame}{Momentos exatos não determinam uma normal}
\begin{columns}[T]
\column{.48\textwidth}\begin{block}{Gerador complexo próprio}\[q_k\sim\CN(0,C_k),\qquad \mathbb E[q_kq_k^T]=0\]\[C_k=P_{\rm GW}(f_k)\Gamma_k+N_k\]\end{block}
\column{.48\textwidth}\begin{block}{Estatísticas quadráticas}\[X_i=q^\dagger H_iq\]\[\mu_i=\operatorname{tr}(H_iC)\]\[\Sigma_{ij}=\operatorname{Re}\operatorname{tr}(H_iCH_jC)\]\end{block}
\end{columns}\take{Mesmo com média e covariância corretas, a distribuição de $X$ pode ter assimetria e caudas não gaussianas.}
\end{frame}
% 11
\begin{frame}{O controle gaussiano isola a compressão}
\begin{columns}[T]
\column{.47\textwidth}\bigword{CN}{\textbf{Observações físicas}}\vspace{3mm}\par Coeficientes complexos geram estimadores quadráticos.\par\vspace{3mm}$A_{\rm CN}$ e $B_{\rm CN}$ usam densidades normais aproximadas.
\column{.47\textwidth}\bigword{G}{\textbf{Controle correspondente}}\vspace{3mm}\par Sorteios normais reais com os mesmos momentos.\par\vspace{3mm}$A_{\rm G}$ e $B_{\rm G}$ têm densidades corretamente especificadas.
\end{columns}\vspace{4mm}\take{No controle G, uma diferença de informação A/B pode ser atribuída à compressão.}
\end{frame}
% 12
\begin{frame}{Congelar a velocidade não é congelar toda a ORF}
\small\renewcommand{\arraystretch}{1.7}\begin{tabular}{lll}\toprule Modelo & Resposta do canal $k$ & O que muda?\\\midrule
$B$ & $\Gamma(\beta_k(u),y_k)$ & Referência dispersiva\\
$C_\beta$ & $\Gamma(\beta_1(u),y_k)$ & Velocidade comum\\
$C_{\rm full}$ & $\Gamma(\beta_1(u),y_1)$ & Resposta inteira do primeiro canal\\\bottomrule\end{tabular}
\vspace{4mm}\take{Todos conservam a dependência da massa. Espectros e ruídos continuam nas frequências físicas.}\source{Definições do experimento tensorial. Na extensão escalar, $C_\beta$ congela também o vínculo escalar.}
\end{frame}
% 13
\begin{frame}{Experimentos complementares, perguntas distintas}
\small\renewcommand{\arraystretch}{1.65}\begin{tabular}{>{\raggedright\arraybackslash}p{3.2cm}>{\raggedright\arraybackslash}p{4.5cm}>{\raggedright\arraybackslash}p{5.3cm}}\toprule Experimento & Desenho & Pergunta\\\midrule
Tensorial de referência & 500 sorteios da priori; 5 modelos & A normal dos estimadores calibra?\\
Geometria ampliada & 12 pulsares/4 canais e 16/8; 596 realizações por configuração & Quanto de informação a compressão perde?\\
Extensão escalar & 1.192 observações; 6.344 análises & Como o vínculo escalar afeta a inferência?\\
Timing e contraste B/C & Projeção temporal; 50 casos condicionais & Quais hipóteses limitam a interpretação?\\\bottomrule\end{tabular}
\end{frame}
% 14
\begin{frame}{A precisão é verificada antes da interpretação}
\begin{columns}[T]
\column{.30\textwidth}\bigword{{\fontsize{20}{24}\selectfont Física}}{Resposta angular}\par\vspace{3mm}Integração no céu e expansão harmônica. Limites analíticos.
\column{.30\textwidth}\bigword{{\fontsize{20}{24}\selectfont Inferência}}{Integração posterior}\par\vspace{3mm}Refinamento, métodos independentes e sensibilidade dos quantis.
\column{.30\textwidth}\bigword{{\fontsize{20}{24}\selectfont Estatística}}{Calibração}\par\vspace{3mm}Realizações da priori, cobertura e propagação do erro numérico.
\end{columns}\vspace{4mm}\take{Concordância numérica e calibração estatística respondem a perguntas diferentes.}
\end{frame}
\section{Resultados: verossimilhança e informação}
% 15
\begin{frame}{Calibração: a verdade ocupa o lugar esperado?}
\centering\begin{tikzpicture}[>=Latex,box/.style={fill=Mist,rounded corners,minimum height=10mm,align=center},node distance=8mm]
\node[box](p){Sortear $\theta_\star$\\da priori};\node[box,right=of p](d){Simular dados\\$d\mid\theta_\star$};\node[box,right=of d](post){Calcular\\$p(\theta\mid d)$};\node[box,right=of post](pit){Avaliar\\$F(\theta_\star\mid d)$};\draw[->,thick](p)--(d);\draw[->,thick](d)--(post);\draw[->,thick](post)--(pit);
\end{tikzpicture}\vspace{5mm}\[\text{PIT}=F(\theta_\star\mid d)\ \sim\ \mathcal U(0,1)\quad\text{sob o modelo correto.}\]
\take{Um intervalo central de 90\% deve conter a verdade em cerca de 90\% das repetições.}\source{Calibração baseada em simulação: Talts et al. (2018); Modrák et al. (2025). Testes com correção de Holm.}
\end{frame}
% 16
\begin{frame}{A cobertura cai para 62,6\% na aproximação A}
\begin{columns}\column{.68\textwidth}\fig{figures/cobertura.pdf}\column{.28\textwidth}\bigword{62,6\%}{Cobertura obtida em $A_{\rm CN}$}\vspace{5mm}\par\small Parâmetro: $\log_{10}\mathrm{EFAC}$.\par\vspace{2mm}500 realizações da priori.\end{columns}
\source{Dissertação, cap. 6. Pontos: coberturas nominais. Barras: sensibilidade à incerteza numérica, não intervalos binomiais. Linha tracejada: 90\%.}
\end{frame}
% 17
\begin{frame}{A distribuição dos PITs revela a inadequação}
\fig[.93]{figures/pit_efac.pdf}
\take{O desvio da diagonal persiste mesmo com os dois primeiros momentos calculados corretamente.}
\source{EFAC, 500 realizações. Faixas coloridas: sensibilidade numérica; tracejado: CDF uniforme. Não são bandas de confiança.}
\end{frame}
% 18
\begin{frame}{Compressão consistente ainda pode perder informação}
\[B=WA,\qquad \mu_B=W\mu_A,\qquad\Sigma_B=W\Sigma_AW^T.\]
\vspace{2mm}\[I(\theta;B)\le I(\theta;A).\]
{\small\textbf{Como quantificar a perda}}\[\left\langle\KL[p(\theta\mid A)\Vert\pi]\right\rangle
-\left\langle\KL[p(\theta\mid B)\Vert\pi]\right\rangle\]
\take{A desigualdade se aplica à média no modelo correto; não ordena toda realização individual.}\source{Desigualdade de processamento de dados; aplicação ao controle gaussiano G. Informação em nats.}
\end{frame}
% 19
\begin{frame}{A geometria de catálogo amplia o experimento}
\begin{columns}\column{.61\textwidth}\fig{../../figures/aplicacao/geometria_academica.pdf}\column{.35\textwidth}\bigword{12 $\to$ 16}{Pulsares}\vspace{4mm}\par\bigword{4 $\to$ 8}{Canais de Fourier}\vspace{4mm}\par\small Posições: catálogo MeerKAT. Distâncias, sinais e ruídos: sintéticos.\end{columns}
\source{Dissertação, cap. 10. Populações aninhadas P12K4 e P16K8; observações pareadas.}
\end{frame}
% 20
\begin{frame}{A perda média é mensurável no controle gaussiano}
\fig[.95]{figures/informacao.pdf}
\take{A informação sobre a massa diminui ao comprimir, nas duas configurações estudadas.}
\source{500 sorteios da priori por configuração; $\langle\KL^{(A)}-\KL^{(B)}\rangle$. Barras: erro padrão Monte Carlo entre realizações; não erro de quadratura. Ruído fixado.}
\end{frame}
% 21
\begin{frame}{A priori sozinha produz um limite superior}
\begin{columns}\column{.56\textwidth}\fig{figures/priori.pdf}\column{.40\textwidth}\[\mathcal L(u)=\text{constante}\]\[\pi(u)=\mathcal U(0,1)\]\[\boxed{Q_{95}=0{,}95,\qquad\KL=0}\]\[m_gc^2\le 0{,}95\frac{h_{\rm P}}{T}\]\end{columns}
\take{Relatar um limite não basta: é preciso mostrar o que a posterior aprendeu em relação à priori.}
\end{frame}
% 22
\begin{frame}{Um limite mais baixo pode vir de um modelo errado}
\begin{columns}[T]\column{.47\textwidth}\begin{block}{Sensibilidade examinada}Medida e suporte da priori.\par\vspace{2mm}Covariância completa ou diagonal.\par\vspace{2mm}Ruído, distâncias e contaminantes.\end{block}\column{.47\textwidth}\begin{block}{Interpretação}Omissões monopolares e dipolares podem deslocar os limites.\par\vspace{3mm}Uma mudança da priori altera a medida de probabilidade atribuída à massa.\end{block}\end{columns}
\take{Robustez é estabilidade sob hipóteses justificadas, não apenas redução do quantil superior.}\source{Dissertação, cap. 8: análises condicionais de robustez.}
\end{frame}
\section{Resultados: frequência, polarização e timing}
% 23
\begin{frame}{Sensibilidade da resposta e o efeito da marginalização}
\begin{columns}[T]\column{.47\textwidth}\bigword{0,01}{Resolução-alvo na largura da priori}\vspace{4mm}\par 24 realizações sorteadas e 8 casos de fronteira.\par\vspace{3mm}Comparação pareada de $\Delta Q_{95}(u)$.\column{.47\textwidth}\begin{block}{Impacto da dimensionalidade}Com ruído marginalizado, os intervalos numéricos se alargam pela integração multidimensional.\par\vspace{3mm}Diferenças pontuais são pequenas, mas exigem isolamento condicional.\end{block}\end{columns}
\take{A alta dimensionalidade mascara contrastes sutis; motivou o teste condicional resolvido a seguir.}\source{Dissertação, cap. 7, tabela 7.5. Erro de integração separado da variabilidade bootstrap.}
\end{frame}
% 24
\begin{frame}{Reduzir a dimensão permite uma pergunta resolvida}
\begin{columns}[T]\column{.30\textwidth}\bigword{50}{Realizações selecionadas antes dos contrastes}\column{.30\textwidth}\bigword{6}{Modelos, nos mesmos dados}\column{.30\textwidth}\bigword{300}{Posteriores condicionais de massa}\end{columns}
\vspace{6mm}Amplitudes e ruído fixados nos valores de geração; somente $u$ é integrado.\par\vspace{4mm}Inversão de envelopes de CDF converte sensibilidade vertical em intervalo do quantil.
\take{É um estudo condicional complementar: a marginalização original continua sendo uma questão distinta.}
\end{frame}
% 25
\begin{frame}{Os deslocamentos médios ficam abaixo da escala-alvo}
\fig[.94]{figures/condicional.pdf}
\take{Todos os 200 contrastes têm raio numérico menor que 0,01; as quatro médias ficam dentro da janela.}
\source{50 realizações P16K8. Pontos: médias; barras: intervalos de sensibilidade numérica das médias. Área clara: $[-0,01;0,01]$. Não são ICs populacionais.}
\end{frame}
% 26
\begin{frame}{No limiar, escalar e tensor têm a mesma forma angular}
\begin{columns}\column{.50\textwidth}\[\beta\to0\quad\Longrightarrow\quad\Gamma_0=\frac12\Gamma_T\]\[S_T\Gamma_T+S_0\Gamma_0=\left(S_T+\frac{S_0}{2}\right)\Gamma_T\]\column{.44\textwidth}\begin{block}{Consequência física}A correlação angular identifica uma combinação de potências.\par\vspace{4mm}Ela não separa $S_T$ e $S_0$ nesse limite.\end{block}\end{columns}
\take{Descoberta analítica: correlações angulares puras no limiar não discriminam modos; exige-se evolução espectral.}\source{Extensão vinculada de Fierz--Pauli, normalização desta pesquisa; dissertação, cap. 9.}
\end{frame}
% 27
\begin{frame}{A componente escalar também exige uma lei adequada}
\begin{columns}\column{.50\textwidth}\bigword{365 / 500}{Cobertura de $u$ em $A_{\rm CN}$: 73\%}\vspace{5mm}\par\bigword{444 / 500}{Cobertura de $u$ em $A_{0,\rm CN}$: 88,8\%}\column{.44\textwidth}\begin{block}{Campanha escalar}Intervalos centrais nominais de 90\%.\par\vspace{3mm}O vínculo entre deformações foi mantido na resposta.\par\vspace{3mm}Ruído e amplitudes auxiliares conhecidos.\end{block}\end{columns}
\take{A falha da normal aplicada a estimadores físicos reaparece na extensão escalar.}\source{Dissertação, cap. 9. Contagens nominais; sensibilidade numérica: [357,369] e [436,444], respectivamente.}
\end{frame}
% 28
\begin{frame}{Ajustar o timing modifica a distribuição de Fourier}
\[r=M\delta\xi+s+n\quad\xrightarrow{\text{ajuste}}\quad r_{\rm pos}=R(s+n)\quad\xrightarrow{\text{Fourier}}\quad z=Dr_{\rm pos}.\]
\vspace{3mm}\begin{columns}[T]\column{.47\textwidth}\begin{block}{Covariância}\[C_z=DRK R^T D^\dagger\]\end{block}\column{.47\textwidth}\begin{block}{Pseudocovariância}\[P_z=DRK R^T D^T\]\end{block}\end{columns}
\take{Ruído gaussiano no tempo não garante coeficientes próprios e independentes depois do ajuste.}\source{$K$: covariância temporal; $R$: projeção do ajuste; $D$: transformação de Fourier.}
\end{frame}
% 29
\begin{frame}{A projeção remove potência nas escalas ajustadas}
\fig[.91]{figures/timing_transfer.pdf}
\source{Experimento sintético de um pulsar: $T=4,5$ anos, 256 tempos; 20.000 séries por amostragem. Rotação: $1,t,t^2$; termos anuais: seno e cosseno. Curvas: resposta da projeção regular.}
\end{frame}
% 30
\begin{frame}{Surgem correlações entre canais e pseudocovariância}
\fig[.92]{figures/timing_matrices.pdf}
\take{No exemplo regular com ajuste anual, a correlação cruzada máxima chega a 0,706.}
\source{Módulos normalizados após ajuste de rotação e termos anuais. O experimento de timing ilustra a hipótese que precisa mudar na aplicação a TOAs.}
\end{frame}
\section{Conclusão}
% 31
\begin{frame}{Roteiro para a aplicação em dados observacionais de TOAs}
\begin{columns}[T]\column{.48\textwidth}\begin{block}{Contribuições estabelecidas}Resposta dispersiva e vínculo de Fierz--Pauli.\par\vspace{0.5mm}Falha de calibração da aproximação normal.\par\vspace{0.5mm}Custo da compressão medido no controle G.\par\vspace{0.5mm}Efeito analítico da priori e quebra por timing.\end{block}\column{.48\textwidth}\begin{block}{Roteiro observacional}Projeção de timing na rede ($G^{\mathsf T}M = 0$).\par\vspace{0.5mm}Tratamento da pseudocovariância ($11{,}22$ nat).\par\vspace{0.5mm}Verossimilhanças de ordem superior.\par\vspace{0.5mm}Dados de banda ultra-larga do SKA.\end{block}\end{columns}
\vspace{-1mm}\take{O trabalho estabelece os critérios diagnósticos e a fundamentação para guiar análises em dados reais.}
\end{frame}
% 32
\begin{frame}{Conclusões centrais e diretrizes desta defesa}
\begin{enumerate}\setlength\itemsep{2mm}
\item \textbf{Momentos corretos não garantem verossimilhança calibrada:} formas quadráticas físicas sofrem quebra de calibração em SBC ($p \le 0{,}005$) e subcobertura nas caudas.
\item \textbf{Custo informacional da compressão quantificado:} o controle gaussiano exato $G$ mede perda média de $0{,}062$ e $0{,}106$ nat nas geometrias MeerKAT.
\item \textbf{Limites superiores exigem métricas de aprendizado:} $Q_{0{,}95} = 0{,}95$ surge puramente da priori sob verossimilhança não-informativa ($D_{\rm KL} = 0$). Relatar $D_{\rm KL}$ é mandatório.
\item \textbf{Degenerescência no limiar e novos horizontes:} $\Gamma_0 = \Gamma_T/2$ no limiar de Fierz--Pauli. O método prepara a transição para dados reais com matriz de timing e dados do SKA.
\end{enumerate}\vspace{2mm}\textcolor{Teal}{\Large Obrigado.}\hfill\small Gráviton de Souza · ITA
\end{frame}
\appendix\backuptrue
\section{Apoio à discussão}
\begin{frame}{Fierz--Pauli linear não encerra a questão não linear}
\begin{itemize}\setlength\itemsep{4mm}\item \textbf{vDVZ:} o modo escalar mantém acoplamento a fontes no limite linear de massa pequena.\item \textbf{Boulware--Deser:} extensões não lineares genéricas perdem um vínculo e podem propagar um modo adicional instável.\item \textbf{Vainshtein:} efeitos não lineares suprimem a contribuição escalar perto de fontes; em modelos apropriados, $r_V\sim(r_s/\mu^2)^{1/3}$.\end{itemize}
\source{Hinterbichler (2012); de Rham (2014); Babichev e Deffayet (2013). A análise PTA usa o regime linear especificado.}
\end{frame}
\begin{frame}{Convenções de base: coeficientes e curvatura}
\[P=I-\widehat k\widehat k,\qquad L=\widehat k\widehat k,\qquad E_0\propto P-2\chi L,\quad\chi=(f_g/f)^2.\]
\small\renewcommand{\arraystretch}{1.7}\begin{tabular}{llll}\toprule Convenção & Base transversal & Base longitudinal & Razão dos coeficientes\\\midrule Bruta & $P$ & $L$ & $-2\chi$\\ Norma dois & $P$ & $\sqrt2L$ & $-\sqrt2\chi$\\ Norma um & $P/\sqrt2$ & $L$ & $-\sqrt2\chi$\\\bottomrule\end{tabular}
\take{As projeções de curvatura usadas no TG satisfazem $p_l/p_b=-\chi$. Um reescalamento relativo da base muda os coeficientes, não a deformação física.}
\end{frame}
\begin{frame}{A forma sinc remove a singularidade aparente}
\[D=1+\beta\widehat\Omega\cdot p,\qquad g(y,D)=\frac{1-e^{-iyD}}{D}\]
\[g(y,D)=iy\,e^{-iyD/2}\operatorname{sinc}(yD/2)\]
\[\operatorname{sinc}(x)=\frac{\sin x}{x},\qquad\operatorname{sinc}(0)=1.\]
\take{O limite $D\to0$ é regular: $g\to iy$. A implementação preserva o termo de pulsar.}
\end{frame}
\begin{frame}{Marginalização analítica da escala comum}
\[a=g-t,\quad b=r-t,\qquad\left|\frac{\partial(g,r,t)}{\partial(a,b,t)}\right|=1\]
\[s=10^t\quad\Longrightarrow\quad dt=\frac{ds}{s\ln10}.\]
\begin{block}{Consequência}A integral sobre a escala comum pode ser reduzida analiticamente, respeitando o suporte transformado das prioris. Na densidade complexa, surgem diferenças de funções gama incompletas.\end{block}
\source{Dissertação, seção 6.3 e apêndice A. O suporte finito e a transformação da medida fazem parte da integral.}
\end{frame}
\begin{frame}{O que significam as barras dos resultados?}
\small\renewcommand{\arraystretch}{1.8}\begin{tabular}{>{\raggedright\arraybackslash}p{4cm}>{\raggedright\arraybackslash}p{8.5cm}}\toprule Figura & Significado\\\midrule Cobertura / PIT & Variação permitida pela incerteza numérica da posterior.\\ Perda média de informação & Erro padrão Monte Carlo entre 500 realizações.\\ Contrastes condicionais & Intervalos numéricos propagados por inversão das CDFs; não ICs de uma população.\\\bottomrule\end{tabular}
\take{Essas três incertezas não são intercambiáveis.}
\end{frame}
\begin{frame}{Referências selecionadas: física}
\small\begin{itemize}\setlength\itemsep{3mm}
\item De Paula, W. L. S. \textit{Estudo da polarização de ondas gravitacionais com gráviton massivo}. TG, ITA, 2003.
\item De Paula, Miranda e Marinho. \textit{Polarization states of gravitational waves with a massive graviton}. CQG, 2004.
\item Liang e Trodden. \textit{Detecting the stochastic gravitational wave background from massive gravity with pulsar timing arrays}. PRD, 2021; versão corrigida.
\item Hinterbichler. \textit{Theoretical aspects of massive gravity}. RMP, 2012.
\item De Rham. \textit{Massive Gravity}. Living Reviews in Relativity, 2014.
\end{itemize}
\end{frame}
\begin{frame}{Referências selecionadas: inferência e materiais}
\small\begin{itemize}\setlength\itemsep{3mm}
\item Talts et al. \textit{Validating Bayesian inference with simulation-based calibration}. arXiv:1805.09294.
\item Modrák et al. \textit{Simulation-Based Calibration Checking for Bayesian Computation: The Choice of Test Quantities Shapes Sensitivity}. Bayesian Analysis, 2025.
\item Franciolini et al. \textit{Likelihoods for Stochastic Gravitational Wave Background Data Analysis}. PRD, 2025.
\item Han e Zhao. \textit{Forecasting graviton-mass constraints from the full covariance of PTA--astrometry ORF estimators}. arXiv:2604.23384.
\end{itemize}\vspace{3mm}\textbf{Dissertação, código e resultados:}\par\href{https://github.com/flavioluiz/TGdoWayne/releases/tag/v1.1.0}{github.com/flavioluiz/TGdoWayne · dissertação v1.1.0}
\end{frame}
\end{document}
